% Kreuzung aus Sicht von oben: «Kein Vortritt» wird mehrfach und damit redundant kodiert
% (Signal, Haifischzähne, Wartelinie). Ersetzt die frühere Rasterillustration.
\begin{tikzpicture}[
        strasse/.style={fill=black!12},
        markierung/.style={white, line width=1.4pt},
        erklaerung/.style={font=\small, align=left, anchor=west},
        pfeil/.style={-{Latex[length=3mm]}, line width=2.5pt}
    ]
    % Vortrittsstrasse und einmündende Strasse
    \fill[strasse] (-5,0) rectangle (5,2.4);
    \fill[strasse] (-1.2,-3.6) rectangle (1.2,0);
    % Mittellinie der Vortrittsstrasse
    \draw[markierung, dashed, dash pattern=on 8pt off 8pt] (-5,1.2) -- (5,1.2);
    % Wartelinie und Haifischzähne (Spitzen zeigen zum wartenden Verkehr)
    \draw[white, line width=2pt, dashed, dash pattern=on 6pt off 4pt] (-1.2,0) -- (1.2,0);
    \foreach \x in {-1.0,-0.6,...,1.0}
        \fill[white] (\x-0.15,-0.1) -- (\x+0.15,-0.1) -- (\x,-0.5) -- cycle;
    % Verkehr
    \node[font=\Large] at (0.6,-2.7) {\faBicycle};
    \draw[pfeil, red!75!black] (0.6,-2.2) -- (0.6,-0.8);
    \node[font=\Large] at (-3.8,0.6) {\faCar};
    \draw[pfeil, green!50!black] (-3.1,0.6) -- (-1.6,0.6);
    % Signal «Kein Vortritt» auf Stange, Spitze nach unten
    \draw[line width=2pt, black!60] (2.1,-0.4) -- (2.1,-1.8);
    \fill[black!60] (2.1,-1.8) circle (2.5pt);
    \draw[red!80!black, line width=3pt, fill=white, line join=round]
        (1.5,0.0) -- (2.7,0.0) -- (2.1,-1.0) -- cycle;
    % Beschriftung
    \node[erklaerung] (a) at (3.4,-0.5) {Signal};
    \node[erklaerung] (b) at (3.4,-1.5) {Haifischzähne};
    \node[erklaerung] (c) at (3.4,-2.5) {Wartelinie};
    \draw[thin, black!70] (a.west) -- (2.75,-0.4);
    \draw[thin, black!70] (b.west) -- (0.9,-0.35);
    \draw[thin, black!70] (c.west) to[out=180, in=-60] (1.15,-0.05);
\end{tikzpicture}
